\documentclass[11pt]{article}
\usepackage[margin=1in]{geometry}
\usepackage{amsmath,amssymb,amsthm,mathtools}
\usepackage[colorlinks=true,linkcolor=blue,citecolor=blue]{hyperref}

\newtheorem{theorem}{Theorem}
\newtheorem{lemma}{Lemma}
\newtheorem{corollary}{Corollary}
\theoremstyle{definition}
\newtheorem{hyp}{Input}
\theoremstyle{remark}
\newtheorem*{remark}{Remark}

\newcommand{\dd}{\,\mathrm{d}}
\newcommand{\R}{\mathbb{R}}
\newcommand{\TW}{\mathrm{TW}}
\newcommand{\Pp}{\mathbb P}

\title{\textbf{Composition error propagation for the noisy folded limit cycle}\\[2pt]
\large the composed chart map concentrates, $\theta$-uniformly: errors add and the fold
renews}
\author{Solomon Williams \\ \small University of Edinburgh}
\date{\today}

\begin{document}
\maketitle

\begin{abstract}\noindent
We prove the propagation step of Path~A: the composed noisy transition map
$\Pi=\Pi_3\circ\kappa_{23}\circ\Pi_2\circ\kappa_{12}\circ\Pi_1$ across the blow-up charts
$K_1$ (entry) $\to K_2$ (inner/fold) $\to K_3$ (exit) concentrates around the
deterministic JKK passage with errors that \emph{add} rather than compound, uniformly in
the angle $\theta\in S^1$. The mechanism is a telescoping (discrete Gronwall) bound for
the transverse deviation together with the \emph{renewal} supplied by the fold: the
entry-to-peel-off map is an exponential contraction, so any entry deviation is forgotten
and cannot accumulate across $K_2$. Consequently the exit law is the inner exit measure
(Tracy--Widom, established separately) transported by the bounded exit chart, plus
sub-dominant Gaussian corrections, and the critical-noise law
$\sigma_*^A=2\sqrt\pi\,\sqrt{\varepsilon_2}\,\sqrt{ac/b}$ survives the composition. The
argument reduces the full composition to per-chart inputs that are either standard
Berglund--Gentz concentration ($K_1,K_3$) or already proved ($K_2$ inner exit measure,
forgetting, $\theta$-uniformity).
\end{abstract}

\section{Framework}

After blow-up the passage is a finite chain of slow--fast SDE segments. We track the
\emph{transverse coordinate} $x$ (the deviation of the noisy trajectory from the
deterministic canard $\gamma$) across a sequence of sections
$\Sigma_0\to\Sigma_1\to\Sigma_2\to\Sigma_3$, with $\Sigma_0$ the entry to $K_1$,
$\Sigma_1=\Sigma_{12}$ the $K_1/K_2$ overlap, $\Sigma_2=\Sigma_{23}$ the $K_2/K_3$
overlap, and $\Sigma_3$ the exit. The chart-overlap maps $\kappa_{i,i+1}$ (smooth
blow-down Jacobians) are absorbed into the transition maps. Writing $x_i$ for the state
on $\Sigma_i$ and $\bar x_i$ for the deterministic state, each stage is a random map
\begin{equation}\label{eq:stage}
  x_i=f_i(x_{i-1})+\zeta_i,\qquad \bar x_i=f_i(\bar x_{i-1}),
\end{equation}
where $f_i$ is the deterministic transition (JKK) and $\zeta_i$ the noise increment of
the $i$th segment. By the strong Markov property the $\zeta_i$ are generated by disjoint
portions of the driving Brownian motion, hence conditionally independent given the
$x_{i-1}$. Set $\delta_i:=x_i-\bar x_i$ (the accumulated transverse error).

We will use two structural inputs that are \emph{proved} elsewhere.

\begin{hyp}[$\theta$-uniform deterministic bounds]\label{in:det}
Each $f_i$ is Lipschitz with constant $L_i$, and the $L_i$ together with the contraction
and quasipotential constants are bounded above and below by $\theta$-independent
constants depending only on $\min/\max$ of $a,b,c$ on $S^1$. \emph{(Proved:
$\theta$-uniformity lemma; the deterministic maps and their derivatives are JKK
Theorem~3.2.)}
\end{hyp}

\begin{hyp}[fold renewal / forgetting]\label{in:forget}
The entry-to-peel-off map (the composition $K_1$ then the attracting phase of $K_2$ up to
the fold) is an exponential contraction: its Lipschitz constant is
$\Lambda_c\le e^{-\frac43 Y_*^{3/2}/1}$ \emph{(}contraction rate $2\sqrt{Y}$ integrated to
the entry level $Y_*$\emph{)}, so $\Lambda_c\to0$. Equivalently $L_2\le\Lambda_c$ in
\eqref{eq:stage}. \emph{(Proved: the canard is attracting; this is the ``forgetting''
Lemma~R1 used in the shooting characterisation, and is confirmed numerically---the
peel-off law is invariant to entry offset/noise within the basin.)}
\end{hyp}

\section{The propagation lemma}

\begin{lemma}[errors add]\label{lem:prop}
For the chain \eqref{eq:stage},
\[
  |\delta_m|\ \le\ \Bigl(\prod_{i=1}^m L_i\Bigr)|\delta_0|
              \ +\ \sum_{j=1}^m\Bigl(\prod_{i=j+1}^m L_i\Bigr)|\zeta_j|.
\]
Moreover, if $\Pp(|\zeta_j|>h_j\mid x_{j-1})\le p_j(h_j)$ for the relevant entry range,
then for any tube radii $h_j$,
\[
  \Pp\Bigl(|\delta_m|>\textstyle\bigl(\prod_i L_i\bigr)|\delta_0|+\sum_j\bigl(\prod_{i>j}L_i\bigr)h_j\Bigr)
  \ \le\ \sum_{j=1}^m p_j(h_j).
\]
All constants are $\theta$-uniform by Input~\ref{in:det}.
\end{lemma}

\begin{proof}
From \eqref{eq:stage}, $\delta_i=x_i-f_i(\bar x_{i-1})
=\zeta_i+\bigl(f_i(x_{i-1})-f_i(\bar x_{i-1})\bigr)$, so by the Lipschitz bound
$|\delta_i|\le|\zeta_i|+L_i|\delta_{i-1}|$. Iterating this scalar recurrence (discrete
Gronwall) gives the deterministic estimate. For the probabilistic bound, the event
$\{|\delta_m|>\cdots\}$ is contained in $\bigcup_j\{|\zeta_j|>h_j\}$ (if every
$|\zeta_j|\le h_j$ the deterministic estimate gives the displayed bound); the union bound
and the conditional tails $p_j$, together with the tower property over the Markov chain,
yield the claim. Uniformity in $\theta$ is inherited from the $\theta$-uniform $L_i$ and
$p_j$ (Input~\ref{in:det}).
\end{proof}

\begin{remark}[why errors do not compound]
The only way the chain could blow up is through $\prod_i L_i\gg1$. Two facts forbid it:
the products are $\theta$-uniformly bounded (Input~\ref{in:det}; JKK's composed map is a
diffeomorphism with controlled derivatives), and---decisively---the fold factor obeys
$L_2\le\Lambda_c\to0$ (Input~\ref{in:forget}). Hence every product that contains the fold
stage is exponentially small. This is the \emph{renewal}: the fold resets the transverse
memory, so deviations entering $K_2$ are annihilated, not amplified.
\end{remark}

\section{Per-chart inputs}

\begin{hyp}[$K_1$, attracting Fenichel tube; Berglund--Gentz]\label{in:K1}
On $K_1$ the deviation from the attracting slow manifold is an Ornstein--Uhlenbeck-type
process with $\theta$-uniform contraction rate $\kappa_1$; for $h>0$,
$\Pp(\sup|\zeta_1|>h)\le C_1\,(t_1/\eta^2)\,e^{-\kappa_1 h^2/2\eta^2}$, and $f_1$ is a
contraction onto $\gamma$. \emph{\cite[Ch.~5]{BG}.}
\end{hyp}

\begin{hyp}[$K_2$, inner exit measure]\label{in:K2}
The peel-off level on $\Sigma_2$ has the inner exit measure $\mu_\eta$: in the inner
scaling $(Y,R)=(\eta^{4/3}y,\eta^{2/3}r)$ it is the law of $-\Lambda_0(\beta)\stackrel{d}{=}
\TW_\beta$, $\beta=4/\eta^2$, \emph{independent of the entry} (Input~\ref{in:forget});
the early-escape probability is $1-F_\beta(0)$. Its tail is the $\TW_\beta$ tail, so
$\Pp(|\zeta_2|>h)\le p_2(h)$ with $p_2$ the (Painlev\'e~II) Tracy--Widom tail.
\emph{(Proved: reduction theorem and shooting characterisation.)}
\end{hyp}

\begin{hyp}[$K_3$, exit tube; Berglund--Gentz]\label{in:K3}
On $K_3$ the trajectory, having left the canard, follows the deterministic exit solution
with a $\theta$-uniform Gaussian tube: $\Pp(\sup|\zeta_3|>h)\le
C_3\,(t_3/\eta^2)\,e^{-\kappa_3 h^2/2\eta^2}$, and the exit map $f_3$ has $\theta$-uniform
Lipschitz constant $L_3\le\bar L$ (the bounded blow-down Jacobian of $\kappa_{23}$ and the
$K_3$ flow). \emph{\cite[Ch.~5]{BG}; \cite{BGK}.}
\end{hyp}

\section{The composition theorem}

\begin{theorem}[composition error propagation]\label{thm:main}
Under Inputs~\ref{in:det}--\ref{in:K3}, the composed noisy transition map
$\Pi=\Pi_3\circ\kappa_{23}\circ\Pi_2\circ\kappa_{12}\circ\Pi_1$ satisfies, with
$\delta_0=0$:
\begin{enumerate}
\item[\textup{(a)}] \textup{(concentration)} for tube radii $h_1,h_2,h_3$,
\[
  \Pp\bigl(|\delta_3|>\bar L\,h_2+h_3+\Lambda_c h_1\bigr)\ \le\ p_1(h_1)+p_2(h_2)+p_3(h_3),
\]
$\theta$-uniformly; the entry contribution $\Lambda_c h_1$ is exponentially small;
\item[\textup{(b)}] \textup{(exit law)} as $\eta\to0$, the exit transverse coordinate
converges in law to $f_3(\mu_\eta)$ ---the inner exit measure transported by the exit
chart--- uniformly in $\theta$; the sub-dominant correction is the $O(\eta)$ Gaussian
$\zeta_3$;
\item[\textup{(c)}] \textup{(critical noise)} the escape probability of the whole passage
equals the inner escape probability $1-F_\beta(0)$ up to the sub-dominant corrections of
\textup{(a)}; hence escape becomes likely at
$\sigma_*^A=2\sqrt\pi\,\sqrt{\varepsilon_2}\,\sqrt{ac/b}$, uniformly in $\theta$.
\end{enumerate}
\end{theorem}

\begin{proof}
Apply Lemma~\ref{lem:prop} with $m=3$ and the per-chart tails of
Inputs~\ref{in:K1}--\ref{in:K3}. The Lipschitz products are
$\prod_{i=1}^3 L_i=L_3 L_2 L_1\le \bar L\,\Lambda_c\,L_1$ and, in the sum,
$\prod_{i>1}L_i=L_3L_2\le\bar L\Lambda_c$ (multiplies $|\zeta_1|$),
$\prod_{i>2}L_i=L_3\le\bar L$ (multiplies $|\zeta_2|$), and the empty product $1$
(multiplies $|\zeta_3|$). With $\delta_0=0$ this gives
$|\delta_3|\le \bar L\Lambda_c\,|\zeta_1|+\bar L\,|\zeta_2|+|\zeta_3|$, and the union
bound of Lemma~\ref{lem:prop} yields (a) (writing $\Lambda_c h_1$ for the exponentially
suppressed entry term).

For (b): take $h_1=O(\eta)$, $h_3=O(\eta)$ (the Berglund--Gentz tube radii at
concentration level), so $\bar L\Lambda_c h_1$ and $h_3$ are $O(\eta)$ or smaller, while
$\zeta_2$ carries the inner scale $O(\eta^{2/3})$ (Input~\ref{in:K2}). Thus
$\delta_3=\bar L\,\zeta_2+O(\eta)$, i.e.\ $x_3=f_3(\bar x_2+\zeta_2)+O(\eta)
=f_3(\mu_\eta)+O(\eta)$ in law, uniformly in $\theta$ (the $O(\eta)$ is the union of the
$\theta$-uniform Gaussian tubes). Because $L_2\le\Lambda_c\to0$, the entry law $x_0$ (and
its fluctuation $\zeta_1$) drops out: the fold renews the chain (Input~\ref{in:forget}),
so the limit is independent of the entry data.

For (c): the passage escapes iff the canard peels off before the fold, an event of the
inner chart with probability $1-F_\beta(0)$ (Input~\ref{in:K2}); by (a) the surrounding
charts perturb this only by the union-bound corrections $p_1+p_3=O(e^{-c'/\eta^2})$, which
are sub-dominant to $1-F_\beta(0)$. The leading-order threshold is unchanged: escape
becomes order-one when the integrated inner hazard $H=\eta^2/4\pi$ is order one, i.e.\ at
$\eta_*=2\sqrt\pi$, giving $\sigma_*^A$. All bounds are $\theta$-uniform.
\end{proof}

\section{Scope}

\paragraph{What is proved here.} The \emph{propagation} itself: that the composed map's
transverse error is the $\theta$-uniformly weighted sum of per-chart fluctuations
(Lemma~\ref{lem:prop}), that the fold renewal annihilates the entry contribution
(so errors do not compound across the non-hyperbolic chart), and that the exit law and
the critical-noise threshold are therefore the inner ones transported by the bounded exit
chart. This is exactly the ``error accumulation, not just per-factor'' step that was the
last structural gap in Part~2.

\paragraph{What is imported.} The per-chart concentration on the \emph{normally
hyperbolic} charts $K_1,K_3$ is standard Berglund--Gentz sample-path concentration
\cite{BG,BGK}; the inner ($K_2$) chart's exit measure, the forgetting, and the
$\theta$-uniformity are proved in the companion notes. The deterministic composed map and
its derivative bounds are JKK Theorem~3.2.

\paragraph{What remains.} (i) The blow-\emph{down} from the blow-up charts to the physical
$(\varepsilon_1,\varepsilon_2)$ variables (the final desingularisation bookkeeping,
turning the $\theta$-uniform chart estimates into a statement for the original system);
(ii) the Channel-B (phase) inner with its bifurcation-specific exponent. Neither affects
the propagation proved here.

\begin{thebibliography}{9}
\bibitem{BG} N.~Berglund, B.~Gentz, \emph{Noise-Induced Phenomena in Slow--Fast
Dynamical Systems: A Sample-Paths Approach}, Springer (2006).
\bibitem{BGK} N.~Berglund, B.~Gentz, C.~Kuehn, \emph{From random Poincar\'e maps to
stochastic mixed-mode-oscillation patterns}, J.\ Dynam.\ Differential Equations
\textbf{27} (2015) 83--136; \emph{Hunting French ducks in a noisy environment},
J.\ Differential Equations \textbf{252} (2012) 4786--4841.
\bibitem{JKK} S.~Jelbart, C.~Kuehn, N.~Kuntz, arXiv:2208.01361 (2024). --- deterministic
backbone, Theorem~3.2, charts $K_1$--$K_3$.
\bibitem{int} Companion notes: \texttt{ShootingCharacterisation\_proof} (inner exit
measure, forgetting), \texttt{InstantonPrinciple\_proof} (rate), and
\texttt{PathA\_endgame\_writeup} (\S\,$\theta$-uniformity).
\end{thebibliography}

\end{document}
